\chapter{Boundaries, Downstream Roles, and the Research Program}\label{ch12:program}

A kernel this small earns its keep in two ways: by what it makes possible
and by what it refuses to permit. The preceding chapters developed the
first; this closing chapter develops the second, and then turns the
refusals into a program. We itemize the excluded transfers --- the pieces
of the wider TKS/CAS research corpus that this artifact deliberately does
not encode, every one of them Ring~4 / \textsc{not licensed} --- examine
in depth the three downstream consumers designed to use the kernel
strictly through its types, walk the full eleven-row proof-obligation
ledger with a substantive discussion of each row, situate the work among
its intellectual neighbors, and conclude.

\section{Excluded transfers}\label{ch12:sec:excluded}

The specification's exclusion list is the artifact's constitution. Each
item below is excluded from the kernel and carries the status
\textsc{open}; as research transfers, all remain Ring~4 /
\textsc{not licensed} --- meaning no artifact, document, or downstream
system may treat the transfer as available, and any future admission
requires its own review, not an appeal to this monograph.

\subsection{Surface syntax}

The kernel neither parses nor evaluates TKS v7.4 or v8 surface syntax. No
tokenizer, grammar, or expression evaluator appears anywhere in the three
modules. This is not a deferred convenience: the production notation
compiler already owns surface syntax, and a second parser in the kernel
would create exactly the dual-authority hazard --- two implementations
with divergent corner cases, each citable as ``the'' semantics --- that
the red-team methodology exists to prevent. The kernel's interface is
typed data (objects, arrows, words over declared generators), constructed
programmatically; text stays outside.

\subsection{The Noetic and CAS mapping}

No identifier in the kernel names a Noetic, a Domain, a Drive, a
Register, or any of the 40 Cognitive States, and no table anywhere maps
canonical constructs to kernel arrows. This is the deepest exclusion, and
the one most tempting to violate informally --- it takes one line of
client code to name an object \texttt{"Mind"} and begin drawing
conclusions. The prohibition is therefore worth stating with its reason:
mapping the canon onto the substrate is a \emph{semantic} act with its
own correctness conditions --- which canonical construct is an object,
which an arrow, which a path; whether canonical composition is kernel
composition --- and none of those conditions have been reviewed. The
mapping is PO-TKS8-CORE-08, \textsc{open / not licensed}; until it
closes, a kernel arrow named after a Noetic is a string, not a claim.

\subsection{The RPM monad}

The kernel contains no RPM data, no monad structure, and no prerequisite
gate. The corrected v7.4 record is precise that RPM is a
section-retraction pair --- one-directional prerequisite satisfaction,
not an adjunction --- and formalizing that structure over kernel types
would require design decisions (what object models a pathway state,
what arrow a satisfaction step) that are exactly the deferred semantic
mapping. Importing the kernel's factorial-power theorem into RPM
reasoning without that mapping would reproduce the invalid power
inference the kernel was built to block.

\subsection{Rewrite systems}

The normal-form machinery of Chapter~8 is denotational: shortlex
minimization over a completed closure of \emph{semantic} objects. It is
not a rewrite system for inherited TKS equations, and no termination or
confluence claim for any such system follows from it. The distinction is
easy to blur because both produce ``canonical forms''; the claim boundary
in the specification is explicit that changing the declared generator
order may change representatives, which no confluent rewriting story
would tolerate. A syntactic rewrite layer, if ever wanted, is a separate
artifact with separate obligations.

\subsection{Noetic Topos 2.0 coverage and modalities}

The companion topos-theoretic development --- sites, coverage, sheaf
conditions, modal operators --- shares a type discipline with the kernel
(Section~\ref{ch12:sec:nt2}) but transfers nothing else. The kernel's
four-atom context lemmas in the Lean seed (Chapter~11) are finite case
analyses housed alongside for toolchain convenience; the presheaf and
sheaf layers themselves are in the Lean open scope and outside the
kernel's Python entirely.

\subsection{Stochastic, quantum, and empirical dynamics}

The deterministic lift is a bookkeeping functor for finite deterministic
maps. Stochastic kernels, quantum channels, contractivity claims, and
any empirical cognitive-dynamics content are excluded --- the permanent
noncontractivity fixture exists precisely to block the most tempting
bridge, the reading of $A_f$ as a channel with norm at most one. A
future stochastic extension would change the mathematics (column-
stochastic real matrices, a genuinely different norm theory) and the
evidence obligations; nothing in this artifact licenses it.

\subsection{The Ten Advances}

No transfer from the Ten Advances manuscript is admitted: none of its
constructions appear in kernel types, none of its claims inherit kernel
evidence. The manuscript remains a research input to the wider program,
Ring~4 / \textsc{not licensed} like every other transfer in this
section.

\begin{boundary}\label{ch12:bdy:excluded}
The exclusions are load-bearing, not provisional politeness. Every
downstream consumer in Section~\ref{ch12:sec:downstream} is designed
so that its use of the kernel would survive an auditor checking each
exclusion above; any design that requires breaching one is, by that
fact, blocked pending the corresponding ledger row.
\end{boundary}

\begin{threat}[Semantic transfer by naming]\label{ch12:thr:naming}
The cheapest breach of the exclusion regime requires no code change at
all: construct kernel objects and arrows whose \emph{presentation
names} are canonical terms --- an object named after a Noetic, a
generator labeled after a Drive --- and then read kernel theorems as
if they were theorems about the named constructs. The kernel's own
design partially blunts this (names are excluded from arrow equality
and carry no semantics anywhere in the code), but no type system can
stop a human from reading meaning into a string. The defense is
procedural and must be stated as such: under the artifact's claim
discipline, a canonical name on a kernel value confers exactly nothing
--- any document deriving a canonical conclusion from a named kernel
computation is making the PO-08 transfer without a license, and the
boundary regression's \texttt{cas\_semantic\_transfer\_status}
assertion is the tripwire that keeps the manifest, at least, honest
about it.
\end{threat}

\begin{threat}[Evidence transfer by association]\label{ch12:thr:evidence}
A subtler breach transfers not semantics but \emph{credibility}: citing
the kernel's proof-and-test discipline as if it attached to a
downstream system that merely uses the kernel. The certified-adapter
design of Section~\ref{ch12:sec:pcba8} shows the correct form ---
certificates bind exact kernel bytes and explicitly carry ``no
automatic CAS semantic theorem'' --- and the incorrect form is any
sentence of the shape ``built on the formally verified TKS kernel,''
which is wrong twice (the kernel is not formally verified, per
Chapter~11; and evidence does not compose across an unverified
boundary). Downstream documents inherit the kernel's obligations, not
its receipts.
\end{threat}

\section{The three downstream consumers}\label{ch12:sec:downstream}

Three companion candidates are designed to consume the kernel strictly
through its public types. None is part of the bound artifact; each is
documented here because the kernel's design choices are partly
explained by the loads these consumers will place on it.

\subsection{CAS Core 8: the fractal bridge}\label{ch12:sec:core8}

CAS Core~8 confronts a legacy: fractal digit sequences (the X.Y.Z
cognitive-fractal addresses of the wider system) whose composition
semantics were historically implicit. Its bridge design uses the kernel
as a type discipline for that legacy. The caller --- not the bridge ---
supplies an explicit catalog of typed arrows, one per legacy digit or
digit transition, declaring for each its source and target object. The
bridge then interprets a digit sequence as a \texttt{SemigroupoidPath}
over the supplied catalog: adjacent-object matching is checked by the
path constructor at every junction, composition is the path composite,
and --- decisively --- \emph{missing digits or type mismatches fail
closed}: an unmapped digit or an ill-typed junction raises rather than
guesses.

Two properties of this design deserve emphasis. First, the bridge is
honest about its authority: as the specification puts it, it ``proves
type-correct composition only, not the meaning of a legacy digit.'' The
caller's catalog is an assertion, not a theorem; the bridge certifies
that \emph{given} the catalog, the sequence composes with exact types.
Whether digit 3 \emph{means} the arrow the caller mapped it to is a
semantic question belonging to PO-08's lane. Second, the choice of
paths rather than endomaps is itself an application of Chapter~6's
discipline: a digit sequence is a heterogeneous generator sequence,
so the bridge exposes composites and never powers --- a legacy
``iterate this address'' idiom would have to produce a certified
endomap first, through the same gate as everyone else.

The failure modes the design forecloses are worth naming, because they
are the failure modes the legacy idiom invites. A digit sequence whose
adjacent stages do not actually connect --- stage output in one state
space, next stage input in another --- historically produced a
computed \emph{something}; under the bridge it produces
\texttt{CompositionTypeError} at the exact junction, with both stage
signatures in the message. A digit with no catalog entry historically
defaulted; under the bridge it fails closed. And a fractal address
treated as ``the same operator at every scale'' --- the reading that
would license powers --- is exactly the invalid power inference of
Chapter~1's audit findings, unwritable here because the bridge's output
type has no power operation. What the bridge cannot do is equally
sharp: it cannot tell a caller that their catalog is \emph{right}. Two
callers with conflicting catalogs for the same legacy corpus will each
get flawlessly typed, mutually contradictory composites. Catalog
review is semantic work, and the bridge's design honestly locates it
outside itself.

\subsection{CAS PCBA 8: the certified-transition adapter}\label{ch12:sec:pcba8}

CAS PCBA~8 (the printed-circuit-board-assembly engineering line)
consumes the kernel from the opposite direction: not legacy addresses
but live component models. Its adapter admits a component transition
function into the kernel only under a certificate that the transition
is \emph{total, finite, deterministic, stateless, and non-stochastic}
--- the five conditions under which a transition genuinely is a
finite-function arrow. A certified transition becomes a kernel arrow;
transition chains become semigroupoid paths; and a feedback loop ---
the case where iteration talk is wanted --- requires a certified
\emph{endomap}, so the factorial-power theorem and the idempotence
machinery become available exactly when the type system says they may.

The adapter's certificate discipline mirrors the manifest's: a
certificate ``binds the exact kernel bytes'' --- it names the hash of
the kernel version against which the adaptation was certified, so a
kernel change invalidates outstanding certificates rather than silently
re-interpreting them. And the specification is explicit that the
certificate ``carries no automatic CAS semantic theorem'': adapting a
solder-paste dispenser's state machine into kernel arrows yields typed
composition guarantees for that machine, not a cognitive claim about
anything. The adapter is thus the template for how empirical
engineering systems touch the kernel: through a narrow, certified,
hash-bound gate, gaining type discipline and gaining nothing else.

Each of the five certificate conditions excludes a specific
misadaptation, and enumerating them shows how much real-world texture
the finite-function abstraction refuses to absorb silently.
\emph{Totality} excludes transitions with undefined regions --- a
component that jams on some input is not a total function until its
jam behavior is modeled as an explicit state. \emph{Finiteness}
excludes continuous state (temperatures, analog positions) until it
is discretized by an explicit, reviewed quantization.
\emph{Determinism} excludes transitions with tolerance bands or race
outcomes --- the single largest gap between datasheet behavior and
function semantics. \emph{Statelessness} excludes hidden memory: a
component whose response depends on history is not one arrow but a
family, and must be certified as a transition on an enlarged explicit
state space. \emph{Non-stochasticity} excludes failure rates and
noise, which belong --- if anywhere --- to the excluded stochastic
extension, not to this lift. A transition failing any condition does
not get a degraded certificate; it gets no certificate, and its chain
gets no kernel typing. The discipline's value to the engineering line
is exactly this refusal: what passes the gate composes with
mathematical guarantees, and what does not pass is visibly,
recordedly outside them.

\subsection{Noetic Topos 2.0: shared type discipline and PO-08}\label{ch12:sec:nt2}

The relationship with Noetic Topos 2.0 is subtler: not consumption of
kernel objects but \emph{convergence on a discipline}. The topos
program admits deterministic maps as morphisms and deliberately
refuses to identify Markov kernels, channels, or enriched sheaves
with functions --- the same determinism boundary the kernel draws
with its lift, arrived at independently from the sheaf-theoretic
side. Both artifacts, that is, encode the decision that the
deterministic/stochastic line is a \emph{type} distinction, not a
parameter. The shared discipline makes the eventual reconciliation
--- matching the kernel's objects and arrows with the topos
development's sites and the corrected v7.4 theorem records --- a
well-posed problem, and that problem has a name: it \emph{is}
PO-TKS8-CORE-08, no more and no less. Until that row closes, the two
artifacts cite each other's boundaries, share a Lean toolchain and
verifier discipline (the four-atom context lemmas of Chapter~11 are
the visible overlap), and transfer nothing semantic.

\section{The obligation ledger, row by row}\label{ch12:sec:ledger}

The eleven-row ledger is reproduced in tabular form in Appendix~B; this
section discusses each row substantively --- what closing it would
require, why it matters, and how it is ordered against the others.

\subsection{PO-TKS8-CORE-01: identity and associativity}

\emph{Status: partial formal verification present.} The category laws
carry a pointwise informal proof, bounded exhaustive tests ($19{,}683$
associativity triples, typed identity families), and exact Lean lemmas
in the \texttt{TypedCategory} module. What remains is not the
mathematics --- the laws are as standard as mathematics gets --- but
the \emph{correspondence}: the Lean lemmas govern a model, and the
equivalence of model and Python is open. Closing the row outright
means either a verified extraction (generate the composition core from
the checked development) or a formalized semantics of the Python
fragment in use; both are heavyweight, and the row is best read as an
anchor whose residual gap is shared with every row that cites Lean
evidence. It matters because every other law in the artifact composes
through this one: functoriality, powering, closure soundness all quote
associativity. Ordering: its residual is subsumed by any serious
attack on PO-07.

\subsection{PO-TKS8-CORE-02: reject every source/target mismatch}

\emph{Status: executable validation.} The rejection behavior is a
finite, directly observable property of four code sites (compose, the
endomap wrapper, the path validator, catalog evaluation), and the
negative suites drive each. The row could be strengthened --- a Lean
model of the \emph{partiality} of composition, or an exhaustive
mismatch sweep over bounded object families --- but its practical
force is already the strongest in the ledger: this is the row that
makes untyped composition unwritable, the first of the three audit
failures the kernel answers. It matters because it is the row
downstream consumers lean on hardest: the fractal bridge's fail-closed
junctions and the PCBA adapter's chain checks are this row wearing
application clothes. No ordering constraints; it is closed at its
evidence level and stays open only in the sense that all executable
validation is bounded.

\subsection{PO-TKS8-CORE-03: powers distinct from paths}

\emph{Status: executable validation.} The separation is enforced by
the runtime type system (no \texttt{power} on paths, endomap
certification at construction) and tested negatively. Its
strengthening would be a formal statement that the exported API admits
no derivation of a path power --- an API-level metatheorem, which is
unusual territory and probably not worth its cost given that the
property is syntactically evident: the method does not exist. The row
matters because it kills the invalid power inference --- the second
audit failure --- and because its discipline is what makes the
factorial theorem \emph{safe} to publish: the theorem is dangerous
precisely when applied to things that are not single endomaps.
Ordering: none; conceptually prior to PO-04.

\subsection{PO-TKS8-CORE-04: factorial idempotence}

\emph{Status: standard result; formalization open.} The mathematics
(Chapter~5) is classical finite-semigroup material \cite{pin}: tails
enter cycles within $|X|-1$ steps, all cycle lengths divide $|X|!$.
The exhaustive 289-endomap check covers carriers through size four.
Closing the row means a Lean proof of the universal statement, which
requires formalized cycle decomposition of functional graphs and
divisibility of factorials --- genuine but well-trodden formalization
work, with likely reusable library support. It matters for two
reasons: it is the theorem downstream feedback-loop analysis will
quote (via the PCBA adapter's certified endomaps), and it is the
simplest ledger row whose closure would demonstrate the seed can
grow beyond definition-adjacent lemmas. Ordering: independent of the
lift rows; a natural first target for seed expansion.

\subsection{PO-TKS8-CORE-05: the exact norm formula}

\emph{Status: informal proof present; universal formalization open.}
The fiber-block Gram proof (Chapter~7) is complete informally, the
structural tests check the formula's consequences on the bounded
family, and Lean checks only the two-state constant expansion.
Closing the row requires formalizing spectral reasoning --- largest
eigenvalue of a PSD Gram matrix, permutation similarity, the all-ones
block spectrum --- or an alternative elementary route (directly
bounding $\lVert A_f v\rVert^2$ by Cauchy--Schwarz per fiber and
exhibiting the maximizing vector), the latter probably being the
formalization-friendly path since it avoids general spectral theory
\cite{horn}. The row matters because the norm formula is the
artifact's main quantitative theorem and the anchor of its
noncontractivity boundary. Ordering: depends on the matrix bridge
(part of PO-11) if done spectrally; the elementary route depends only
on arbitrary-vector linearity.

\subsection{PO-TKS8-CORE-06: shortlex soundness and leastness}

\emph{Status: informal correctness argument; formalization open.} The
pruning lemma's induction (Chapter~8) and the finiteness argument are
short but subtle --- the suffix-transfer step is exactly the kind of
argument where informal proofs hide gaps. Closing the row means
formalizing the closure algorithm's invariants: heap order delivers
shortlex order, first insertion is least, pruning preserves
reachability. This is program-adjacent verification (an algorithm,
not just a definition), so it likely proceeds by modeling the
algorithm abstractly in Lean rather than verifying the Python.
It matters because canonical forms are only as trustworthy as
leastness: a wrong representative is not detectably wrong to a
consumer, unlike a type error. Ordering: independent of the lift
rows; benefits from whatever induction infrastructure PO-04 builds.

\subsection{PO-TKS8-CORE-07: independent theorem-prover formalization}

\emph{Status: partial; present.} Chapter~11 is this row in full: the
pinned seed exists with its axiom audit, and the open scope is
itemized. Closing it means discharging the open-scope items --- the
universal norm, the matrix bridge, finite-sum multiplication,
arbitrary-vector linearity, factorial power, shortlex leastness ---
at which point rows 01, 04, 05, 06, and 11 all upgrade
simultaneously. The row matters as the aggregate index of formal
depth, and its ordering is the reverse of the others': it closes
last, when its constituents do. A second-prover corroboration (Coq
or Agda, both \textsc{unavailable} in the probe environment) would
strengthen it further but is not, on the ledger's text, required.

\subsection{PO-TKS8-CORE-08: reconcile kernel types with corrected v7.4 records}

\emph{Status: open; not licensed.} The deepest row, and the only one
whose closure is a semantic act rather than a mathematical or
engineering one. Closing it requires: a reviewed mapping proposal
(which canonical constructs correspond to objects, arrows, endomaps,
paths); verification that the corrected v7.4 theorem records ---
admissibility relations, the section-retraction RPM structure, the
open-simplex state space --- are representable, or explicitly not
representable, in kernel terms; and a licensing decision by review,
not by drafting. It matters more than any other row because it is
the bridge over which the kernel would stop being ``an independent
reference model'' and start carrying canonical content; every
downstream consumer's semantic silence is provisional on this row.
Ordering constraints are the strictest in the ledger: it should not
close before PO-10 (review) has structure, and nothing in Ring~4
crosses until it does.

\subsection{PO-TKS8-CORE-09: integrate with the active notation compiler}

\emph{Status: open (outside this lane).} The kernel and the
production compiler currently share nothing but intent. Integration
--- whether as a semantic checking layer behind the compiler's kinds,
as a validation oracle in compiler tests, or as a type-annotation
target --- requires a compatibility design: mapping compiler kinds
onto kernel types, deciding failure behavior when the kernel rejects
what the compiler accepts, and versioning the joint artifact. The
manifest's integration boundary (\texttt{active\_cas\_compiler\_modified:
false}, \texttt{production\_runtime\_dependency: false}) is the
row's current truth, and the boundary regression of Chapter~10 tests
it. The row matters because the audit findings that motivated the
kernel are findings about the production system; the kernel repays
its cost only when its discipline reaches production. Ordering: after
PO-08 for any semantically loaded integration; a purely structural
integration (type-checking composition of compiler programs) could
in principle proceed earlier under its own review.

\subsection{PO-TKS8-CORE-10: independent specialist review}

\emph{Status: open; no external review receipt.} Every evidence
channel in the artifact --- proofs, tests, seed --- is
project-authored. Closing this row means an external reviewer with
category-theoretic and formal-methods competence examines the claims,
the code, and the seed, and issues a receipt that the artifact then
binds. It matters because it is the only row that addresses
statement-selection risk (Chapter~11's closing observation): insiders
choose what to prove and what to test, and only outside eyes audit
the choosing. Ordering: it gates PO-08's closure and should precede
any licensing decision; it can begin at any time --- nothing in the
ledger blocks review, and review would sharpen everything else.

\subsection{PO-TKS8-CORE-11: lift preserves identities and composition}

\emph{Status: partial formal verification; finite-sum matrix bridge
open.} The functor laws hold on $1{,}678$ exact bounded pairs and in
Lean at entry, basis, and support levels; what is missing is the
finite-sum matrix multiplication bridge that would let the
basis-level Lean statements imply the entrywise matrix statements the
Python tests check. Closing it is the most mechanical of the open
formalizations --- finite sums over Kronecker deltas --- and would
knit the two evidence channels together at their closest point of
approach. It matters because functoriality is what makes the lift a
\emph{functor} rather than a family of matrices, and because its
closure is a prerequisite for the spectral route to PO-05. Ordering:
before PO-05 (spectral route); independent otherwise.

\subsection{Ordering constraints in one view}

Collecting the per-row ordering remarks into a single dependency
sketch:

\begin{center}
\small
\begin{tabular}{@{}lp{8.6cm}@{}}
\toprule
Row & Depends on / naturally follows \\
\midrule
PO-01 residual & subsumed by PO-07's open scope (model equivalence apart) \\
PO-02, PO-03 & no dependencies; closed at their evidence level \\
PO-04 & independent; natural first seed-expansion target \\
PO-05 & elementary route: linearity item of PO-07 scope; spectral route: PO-11 bridge \\
PO-06 & independent; shares induction infrastructure with PO-04 \\
PO-07 & aggregates: closes when 04, 05, 06, 11 formalizations land \\
PO-08 & gated by PO-10; gates semantic use everywhere \\
PO-09 & semantically loaded form gated by PO-08; structural form separately reviewable \\
PO-10 & no dependencies; can begin immediately; gates PO-08 \\
PO-11 & bridge is mechanical; precedes spectral PO-05 \\
\bottomrule
\end{tabular}
\end{center}

The sketch exposes the program's critical path: PO-10 (review) is the
only row that gates the semantic future and depends on nothing, which
makes commissioning review the highest-leverage single action
available; and the formalization rows form an almost-parallel bundle
whose one internal edge (bridge before spectral norm) dissolves if
the elementary norm route is taken. A program that started review and
the two mechanical Lean items simultaneously would be attacking every
hub at once.

\subsection{Reading the ledger whole}

Three structural observations. First, the ledger separates
\emph{evidence upgrades} (04, 05, 06, 11, and ultimately 07: open
formalizations of proved mathematics) from \emph{semantic gates} (08,
09: acts requiring review and licensing) from \emph{process
obligations} (10: review itself); the three kinds close by different
means and no amount of progress in one kind substitutes for another.
Second, the dependency structure is shallow --- most rows are
independent --- with two hubs: PO-07 aggregates the formalization
rows, and PO-10 gates the semantic rows. Third, the ledger is
conserved under this monograph: drafting these chapters upgraded no
row, and the artifact's status line --- \textsc{gate 1 prototype /
ring 4 / not licensed} --- is exactly what it was when the manifest
was written. The program is the ledger; the monograph is its
documentation.

\section{Related work}\label{ch12:sec:related}

The mathematics deployed here is deliberately standard. The categorical
framework --- objects, arrows, functors, the discipline of typed
composition --- is textbook material \cite{maclanecwm}, and the kernel's
category is among the simplest nontrivial examples in that literature: a
skeleton-free category of finite sets with named objects. The
finite-semigroup facts underlying the power discipline, including the
idempotent-power phenomenon in finite semigroups, belong to the
tradition surveyed by Pin \cite{pin}. The spectral analysis behind the
norm theorem --- Gram matrices, permutation similarity, the all-ones
block spectrum --- is standard matrix analysis \cite{horn}. The
artifact's registered position is that none of its mathematical content
is new: the theorem ledger records standard results and definitions,
and no novelty search has been commissioned.

The methodological neighbors are likewise explicit. Bounded-exhaustive
checking as a design philosophy descends from the small-scope
hypothesis and the Alloy tradition of Jackson \cite{jackson};
Chapter~10 details both the debt and the contrast with sampled
property-based testing in the QuickCheck lineage \cite{quickcheck}.
The formal seed follows standard kernel-checking practice in Lean~4
\cite{lean4}, with the axiom-audit discipline --- single admissible
standard axiom, mechanical per-theorem dependency reports --- as its
distinguishing procedural feature.

What the artifact offers as a contribution is consequently not
mathematics but \emph{engineering discipline}: name-inclusive object
identity, ordered-input rejection of hash-order nondeterminism, the
endomap/path type separation, fail-closed closure semantics,
hash-bound artifacts with an explicit self-binding boundary, and a
claim ledger that tracks evidence kind per statement. Each element is
individually known practice somewhere; their joint application to a
cognitive-architecture semantic substrate, with the claim discipline
held as strictly as the proofs, is the design.

Three contrasts locate the design more precisely among its neighbors.
Against \emph{model finders} in the Alloy tradition \cite{jackson},
the kernel inverts the workflow: Alloy checks declarative models
against small scopes to find counterexamples before implementation,
whereas this artifact implements first and then exhausts small scopes
to tie the implementation to prior proofs --- the small-scope
hypothesis serving as post-hoc validation warrant rather than
design-time exploration. Against \emph{property-based testing
practice} \cite{quickcheck}, the kernel's suites give up generator
generality (there is no combinator library, no shrinking) in exchange
for countable, derivable, deterministic coverage --- a trade sensible
only because the instance spaces are tiny, and explicitly the wrong
trade for the large-state systems where the QuickCheck tradition
thrives. Against \emph{verified-kernel developments} in the Lean
ecosystem \cite{lean4}, the artifact is deliberately more modest ---
a seed with an audited boundary rather than a formalized theory ---
and correspondingly more explicit than is customary about the
unformalized remainder, which the ledger enumerates item by item
rather than leaving implicit in what the repository happens not to
contain. The standard mathematical content, meanwhile, is simply
inherited with attribution --- category theory \cite{maclanecwm},
finite semigroups \cite{pin}, matrix analysis \cite{horn} --- with
the ledger's \textsc{standard result} tag marking exactly where the
artifact claims curation rather than creation.

\section{Conclusion}\label{ch12:sec:conclusion}

TKS v8.0, as documented across this monograph, is a small object with
sharp edges. Its positive content fits in a sentence each: typed
composition that cannot silently mistype; powers that exist only where
a monoid exists, with the factorial bound as the one safe general
theorem; a deterministic lift whose exact norm is proved rather than
assumed, permanently accompanied by the counterexample that blocks its
over-reading; and canonical forms that are deterministic because their
inputs are ordered and their closure completes or refuses. Its
evidence is layered and labeled: informal proofs for everything,
bounded-exhaustive executable validation ($19{,}683$ triples, $289$
endomaps, $499$ typed maps, $1{,}678$ lift pairs, 20 tests green at
the 2026-08-04 snapshot), and a pinned, axiom-audited Lean seed for
the definition-adjacent core --- with the gaps between layers
itemized, not blurred.

Its negative content is the part most easily lost in summary and most
important to keep: the kernel encodes no Noetic, no Drive, no
Register; it parses nothing; it claims nothing empirical; its
manifest is not yet outer-bound; its suite proves nothing beyond its
bounds; its seed verifies a model, not the program; and every
transfer that would make it mean something canonically remains
Ring~4 / \textsc{not licensed} behind an eleven-row ledger whose
open rows close by proof, by engineering, or by review --- never by
prose. The three downstream consumers show that this austerity is
usable: legacy addresses, certified component transitions, and a
companion topos program can each take exactly the discipline they
need through the types, and nothing more, today.

The research program, then, is not a promise appended to a result;
it is the ledger read forward. Formalize the factorial theorem and
the shortlex invariants; build the finite-sum bridge and take the
elementary route to the universal norm; commission the review that
audits the choosing; and only then, with receipts in hand, open the
question the whole design defers --- whether and how the canon maps
onto the substrate. The kernel's refusal to answer that question
prematurely is not a limitation of the work. It is the work.
